
\documentclass[11pt]{article}
\usepackage[margin=1in]{geometry}
\usepackage[T1]{fontenc}
\usepackage{hyperref}
\usepackage{longtable}
\usepackage{listings}

\title{ROS 2 Architecture and Development Guide}
\author{}
\date{\today}

\begin{document}

\maketitle
\tableofcontents
\newpage

\section{Introduction}

ROS 2 (Robot Operating System 2) is a distributed middleware framework used to build robotic systems.
ROS 2 is not an operating system. It is a collection of libraries, tools, communication middleware,
build systems, and conventions that help developers build modular robotic applications.

Major goals of ROS 2:

\begin{itemize}
\item Distributed communication
\item Real-time support
\item Multi-platform support
\item Security
\item Scalability
\item DDS-based communication
\end{itemize}

\section{ROS 2 Architectural Layers}

\subsection{Application Layer}

Contains user-developed nodes.

Examples:

\begin{itemize}
\item Sensor drivers
\item Perception nodes
\item Navigation nodes
\item Control nodes
\item Visualization nodes
\end{itemize}

\subsection{Client Libraries}

Provide APIs used by developers.

Common libraries:

\begin{itemize}
\item rclcpp (C++)
\item rclpy (Python)
\end{itemize}

Responsibilities:

\begin{itemize}
\item Node creation
\item Publisher creation
\item Subscriber creation
\item Services
\item Actions
\item Parameters
\end{itemize}

\subsection{RCL Layer}

The ROS Client Library (RCL) provides a language-independent abstraction.

Responsibilities:

\begin{itemize}
\item Graph management
\item Discovery support
\item Topic management
\item Service management
\item Parameter handling
\end{itemize}

\subsection{RMW Layer}

RMW (ROS Middleware Interface) abstracts DDS implementations.

Examples:

\begin{itemize}
\item Fast DDS
\item Cyclone DDS
\item Connext DDS
\end{itemize}

The application code remains unchanged regardless of DDS vendor.

\subsection{DDS Layer}

DDS (Data Distribution Service) provides:

\begin{itemize}
\item Discovery
\item Publish-subscribe communication
\item Quality of Service (QoS)
\item Reliable transport
\item Best-effort transport
\end{itemize}

\section{ROS 2 Graph}

A ROS graph consists of:

\begin{itemize}
\item Nodes
\item Topics
\item Services
\item Actions
\item Parameters
\end{itemize}

Example:

\begin{verbatim}
Talker Node ---> Topic (/chatter) ---> Listener Node
\end{verbatim}

\section{Nodes}

A node is an executable participating in the ROS graph.

Responsibilities:

\begin{itemize}
\item Publish messages
\item Subscribe to messages
\item Provide services
\item Consume services
\item Manage parameters
\end{itemize}

\section{Topics}

Topics implement asynchronous communication.

Characteristics:

\begin{itemize}
\item One-to-many communication
\item Many-to-many communication
\item Decoupled publishers and subscribers
\end{itemize}

\section{Services}

Services implement request-response communication.

Example:

\begin{verbatim}
Client ---> Request ---> Server
Client <--- Response --- Server
\end{verbatim}

\section{Actions}

Actions are used for long-running tasks.

Example:

\begin{itemize}
\item Navigation
\item Manipulator motion
\item Autonomous missions
\end{itemize}

Actions support:

\begin{itemize}
\item Feedback
\item Cancellation
\item Progress monitoring
\end{itemize}

\section{ROS 2 Development Workflow}

\subsection{Workspace}

A workspace contains source code and build artifacts.

Typical structure:

\begin{verbatim}
ros2_ws/
 ├── src/
 ├── build/
 ├── install/
 └── log/
\end{verbatim}

\subsection{Creating a Workspace}

\begin{verbatim}
mkdir -p ~/ros2_ws/src
cd ~/ros2_ws
colcon build
\end{verbatim}

\subsection{Building}

ROS 2 uses Colcon.

\begin{verbatim}
colcon build
\end{verbatim}

\subsection{Sourcing}

\begin{verbatim}
source install/setup.bash
\end{verbatim}

\section{Talker and Listener Deep Dive}

The demo consists of:

\begin{itemize}
\item demo\_nodes\_cpp talker
\item demo\_nodes\_py listener
\end{itemize}

Topic:

\begin{verbatim}
/chatter
\end{verbatim}

Message Type:

\begin{verbatim}
std_msgs/msg/String
\end{verbatim}

\section{Process Creation}

When you execute:

\begin{verbatim}
ros2 run demo_nodes_cpp talker
\end{verbatim}

The following happens:

\begin{enumerate}
\item Shell launches process.
\item Linux kernel creates process.
\item Executable is loaded into memory.
\item Shared libraries are loaded.
\item main() begins execution.
\item rclcpp::init() initializes ROS.
\item DDS participant is created.
\item Node object is created.
\item Publisher is created.
\end{enumerate}

\section{Talker Startup Sequence}

Internal flow:

\begin{verbatim}
main()
 |
 +--> rclcpp::init()
 |
 +--> create node
 |
 +--> create publisher
 |
 +--> create timer
 |
 +--> spin()
\end{verbatim}

Publisher:

\begin{verbatim}
publisher_ =
create_publisher<std_msgs::msg::String>(
    "chatter", 10);
\end{verbatim}

Timer fires every 500 ms.

\section{Message Creation}

Each timer callback:

\begin{enumerate}
\item Creates String message.
\item Fills text.
\item Increments counter.
\item Calls publish().
\end{enumerate}

Example:

\begin{verbatim}
Hello World: 42
\end{verbatim}

\section{Serialization}

The publisher does not directly send a C++ object.

Steps:

\begin{enumerate}
\item Message created in memory.
\item ROS middleware serializes data.
\item DDS packet generated.
\item Network transport selected.
\end{enumerate}

\section{DDS Discovery}

Before communication begins:

\begin{enumerate}
\item DDS participant announces itself.
\item Listener announces itself.
\item Discovery protocol exchanges metadata.
\item Topic compatibility verified.
\item Communication path established.
\end{enumerate}

No ROS master is required in ROS 2.

\section{Listener Startup}

Command:

\begin{verbatim}
ros2 run demo_nodes_py listener
\end{verbatim}

Startup flow:

\begin{verbatim}
main()
 |
 +--> rclpy.init()
 |
 +--> create node
 |
 +--> create subscriber
 |
 +--> spin()
\end{verbatim}

Subscriber registers callback.

\section{Message Reception}

Sequence:

\begin{enumerate}
\item DDS receives packet.
\item Packet deserialized.
\item Message reconstructed.
\item Callback scheduled.
\item Executor invokes callback.
\item Message printed.
\end{enumerate}

Callback:

\begin{verbatim}
def listener_callback(msg):
    print(msg.data)
\end{verbatim}

\section{End-to-End Message Flow}

\begin{verbatim}
Application (Talker)
      |
      v
rclcpp
      |
      v
RCL
      |
      v
RMW
      |
      v
DDS Writer
      |
      v
Network / Shared Memory
      |
      v
DDS Reader
      |
      v
RMW
      |
      v
RCL
      |
      v
rclpy
      |
      v
Application (Listener)
\end{verbatim}

\section{Executors}

Executors process callbacks.

Types:

\begin{itemize}
\item SingleThreadedExecutor
\item MultiThreadedExecutor
\end{itemize}

Responsibilities:

\begin{itemize}
\item Timer callbacks
\item Subscriber callbacks
\item Service callbacks
\item Action callbacks
\end{itemize}

\section{Quality of Service}

QoS controls communication behavior.

Important settings:

\begin{itemize}
\item Reliability
\item Durability
\item History
\item Deadline
\item Lifespan
\end{itemize}

Common reliability modes:

\begin{itemize}
\item Reliable
\item Best Effort
\end{itemize}

\section{Memory and Threading}

Typical talker execution:

\begin{itemize}
\item Main thread
\item Executor loop
\item DDS background threads
\item Network threads
\end{itemize}

\section{Debugging Tools}

Useful commands:

\begin{verbatim}
ros2 node list
ros2 topic list
ros2 topic echo /chatter
ros2 topic info /chatter
ros2 doctor
\end{verbatim}

\section{Recommended Learning Path}

\begin{enumerate}
\item Nodes
\item Topics
\item Services
\item Actions
\item Parameters
\item Launch files
\item Custom messages
\item Lifecycle nodes
\item Navigation stack
\item MoveIt
\item Real-time systems
\end{enumerate}

\section{Summary}

ROS 2 applications are composed of nodes communicating through DDS.
The talker creates messages, DDS discovers subscribers, messages are serialized,
transported, deserialized, and delivered through executors to subscriber callbacks.
Understanding this end-to-end flow is fundamental for designing scalable robotic systems.

\end{document}
